\clearpage
\pagestyle{empty}
\begin{vietnamese}
{\centering
\thispagestyle{empty}
\linespread{1.15}\fontsize{16pt}{22pt}\selectfont\textbf{CITRA: NỀN TẢNG NƠ-RON -- KÝ HIỆU PHỤC VỤ\\CHẤM ĐIỂM TỰ ĐỘNG TÍNH TOÀN VẸN TRÍCH DẪN\\VÀ PHÁT HIỆN ẢO GIÁC NGUỒN TÀI LIỆU}\par}

\vspace{0.35\baselineskip}

{\centering\linespread{1.15}\fontsize{16pt}{22pt}\selectfont\textbf{TÓM TẮT}\par}

\vspace{0.45\baselineskip}
\thispagestyle{empty}

Các mô hình ngôn ngữ tạo sinh giúp tăng tốc soạn thảo học thuật nhưng thường sinh ra trích dẫn đúng cú pháp mà trỏ tới tài liệu không tồn tại, bài báo đã bị rút lại hoặc sai lệch siêu dữ liệu. Các công cụ kiểm tra đạo văn và văn bản máy sinh hiện có chỉ xét bề mặt câu chữ mà bỏ qua cơ sở dữ liệu đối chứng, trong khi các mô hình học máy thuần túy thất bại trước vi phạm nguồn gốc rời rạc và nhầm lỗi nghẽn mạng với tài liệu bịa đặt. Để lấp khoảng trống này, đồ án tốt nghiệp đề xuất \textbf{CITRA}, nền tảng kiểm chứng trích dẫn theo kiến trúc lai Nơ-ron -- Ký hiệu, thực hiện liên kết hai chiều trên 12 chuẩn trích dẫn và truy hồi liên hợp qua kho cục bộ 129,137 bản ghi (128,115 bài báo duy nhất đủ chuẩn đo kiểm) cùng 15+ cơ sở dữ liệu học thuật và kỹ thuật quốc tế. Hệ thống dung hợp đối sánh ngữ nghĩa song ngữ, cây luật ký hiệu phân cấp ưu tiên, hồi quy logistic đa lớp và hiệu chỉnh xác suất Platt kèm cơ chế từ chối phán quyết có chọn lọc. Đánh giá trên tập kiểm thử tầng quyết định độc lập 18,000 trích dẫn (từ bộ chuẩn 120,000 mẫu), bộ kiểm thử ngoài phân phối 2,400 mẫu, 120 tài liệu tham khảo kiểm chứng trực tuyến đa cơ sở dữ liệu và 310 bản thảo đa định dạng cho thấy CITRA đạt độ chính xác 99.67\% và Macro-F1 99.65\% (giảm 85.7\% tổng lỗi so với mô hình học máy mạnh nhất, $p < 10^{-72}$, 0 kết tội nhầm trên 6,480 trích dẫn hợp lệ), giảm 83.6\% sai số hiệu chỉnh kỳ vọng (ECE) xuống 0.0405, nâng độ chính xác ngoài phân phối từ 57.33\% lên 99.96\%, đạt 89.17\% độ chính xác trực tuyến 5 lớp (F1 lớp hợp lệ 98.36\%), và duy trì độ phủ trích xuất 97.7\%--100\% với độ ổn định xếp hạng Điểm toàn vẹn trích dẫn (CIS) Spearman $\rho = 0.9987$.

\vspace{0.35\baselineskip}
\noindent\textbf{Từ khóa:} Kiểm chứng tính toàn vẹn trích dẫn, phát hiện ảo giác nguồn tài liệu, trí tuệ nhân tạo Nơ-ron -- Ký hiệu, từ chối phán quyết có chọn lọc, hiệu chỉnh xác suất Platt, kỹ thuật phần mềm thực nghiệm.
\end{vietnamese}
\clearpage
